\chapter*{Colophon}
\addcontentsline{toc}{chapter}{Colophon}
\markboth{Colophon}{Colophon}

This manuscript was written in Indonesian, arranged into six parts and thirty-two
chapters, and built with \LaTeX{} (pdf\LaTeX{}, the \texttt{book} class). The entire
manuscript source in \LaTeX{} form ships with this book, so that every number and every
code fragment can be checked again by the reader; the thirteen-chapter version that
preceded it remains archived as a record of how the book was put together.

This book is an original piece of writing: its sentences, its order of presentation, its
worked examples, and its exercises were composed by the author. Its mathematical ideas
come from the books in the References, and the reader is urged to read them. No part of
the manuscript is copied verbatim from a source book; where a definition or a theorem
follows a source closely, that source is named in the chapter map.

Every number inside a computing Lab comes from a real execution on the author's machine
(Python 3, NumPy, SciPy) and can be repeated by the reader. The Lean~4 code fragments in
Chapter~28 were run with Lean 4.10.0 without Mathlib and their output is quoted as it
stood; library tactics such as \texttt{linarith} and \texttt{ring} were not run because
Mathlib is not installed, and that chapter states the boundary.

Cover and layout are deliberately plain, without photographs or raster images: every
figure inside the chapters is drawn as TikZ code, so that the reader can examine its
numbers and change them, and the list of figures is in the front matter. Seven study boxes
serve as the main visual markers --- \emph{For the young reader}, \emph{Note for the
instructor}, \emph{For the researcher}, \emph{Read aloud}, \emph{Work it out by hand},
\emph{Step recipe}, and \emph{Challenge} --- so that one and the same chapter can be
entered through different ways of learning.

\subsection*{License}
The text of this book is licensed under Creative Commons BY 4.0 (CC BY 4.0): it may be
shared, adapted, and used commercially as long as the author is credited and changes are
marked. The \texttt{LICENSE} and \texttt{CITATION.cff} files in the repository carry the
same terms. The Python and Lean~4 code fragments inside it are licensed under MIT, so that
they can be reused without friction.

\subsection*{How to cite}
Pangimpian, Afiq Andico. \emph{Why Machines Can Be Wrong: From Linear Algebra to
Machine-Learning Security}. First Edition. 2026.

\subsection*{Version and feedback}
Manuscript version 1.0, October 2026. Every error you find --- a wrong computation, a
reference that does not match, or a misleading sentence --- is more useful than praise;
send it to \texttt{afiqandico13@gmail.com} with the chapter and page number.
